\documentclass[11pt]{article}

% ---------- PACKAGES ----------
\usepackage[margin=1in]{geometry}
\usepackage{enumitem}
\usepackage{titlesec}
\usepackage{xcolor}
\usepackage{tcolorbox}
\tcbuselibrary{skins,breakable}
\usepackage{hyperref}
\usepackage{longtable}
\usepackage{array}
\usepackage{booktabs}
\usepackage{fancyhdr}
\usepackage{amsmath,amssymb}
\usepackage{parskip}

% ---------- COLORS ----------
\definecolor{jhublue}{RGB}{0,45,114}
\definecolor{jhulight}{RGB}{232,238,247}
\definecolor{diffred}{RGB}{155,28,28}
\definecolor{diffbg}{RGB}{253,235,235}
\definecolor{prereqgreen}{RGB}{31,94,58}
\definecolor{prereqbg}{RGB}{232,245,237}

% ---------- SECTION STYLING ----------
\titleformat{\section}
  {\normalfont\Large\bfseries\color{jhublue}}
  {\thesection}{0.6em}{}
\titleformat{\subsection}
  {\normalfont\large\bfseries\color{jhublue}}
  {\thesubsection}{0.6em}{}
\titlespacing*{\section}{0pt}{18pt}{8pt}
\titlespacing*{\subsection}{0pt}{14pt}{6pt}

% ---------- HEADER/FOOTER ----------
\pagestyle{fancy}
\fancyhf{}
\renewcommand{\headrulewidth}{0.4pt}
\lhead{\small AS.110.106 --- Calculus I (Bio.\ \& Soc.\ Sci.)}
\rhead{\small Course Breakdown}
\cfoot{\thepage}

% ---------- CUSTOM BOXES ----------
\newtcolorbox{modulebox}[1]{
  colback=jhulight, colframe=jhublue, coltitle=white,
  fonttitle=\bfseries, breakable, boxrule=0.8pt, arc=2pt,
  title={#1}
}
\newtcolorbox{diffbox}[1][]{
  colback=diffbg, colframe=diffred, coltitle=white,
  fonttitle=\bfseries\small, breakable, boxrule=0.6pt, arc=1pt,
  title=Difficult / High-Friction Topics, left=6pt, right=6pt, top=4pt, bottom=4pt
}
\newtcolorbox{prereqbox}[1][]{
  colback=prereqbg, colframe=prereqgreen, coltitle=white,
  fonttitle=\bfseries\small, breakable, boxrule=0.6pt, arc=1pt,
  title=Prerequisite Knowledge, left=6pt, right=6pt, top=4pt, bottom=4pt
}

\newcommand{\loheader}{\textbf{Learning Objectives.} By the end of this module, students will be able to:}

% ---------- TITLE ----------
\title{\vspace{-1.5cm}\textcolor{jhublue}{\Huge\bfseries AS.110.106 Calculus I}\\[4pt]
\Large (Biology \& Social Sciences)\\[2pt]
\large Full Course Breakdown: Structure, Modules, Objectives, and Difficulty Map}
\author{Department of Mathematics \\ Johns Hopkins University}
\date{Text: \emph{Calculus for Biology and Medicine}, 4th ed., Neuhauser \& Roper \\ Pearson, 2018 --- ISBN-13: 978-0134070049}

\begin{document}

\maketitle
\thispagestyle{empty}

\begin{abstract}
\noindent This document reorganizes the official AS.110.106 syllabus into a structured
course map: six sequential modules spanning approximately 12--14 weeks, each broken into
chapters/sections, core concepts, explicit learning objectives, historically difficult
topics, and the prerequisite knowledge each module assumes. It is intended as a planning
and study reference that mirrors the syllabus's official topic list while adding
pedagogical structure.
\end{abstract}

\tableofcontents
\vspace{0.5cm}

% =========================================================
\section{Course Structure at a Glance}
% =========================================================

The course is organized into six sequential modules that build on one another. Total
pacing is approximately 12--14 weeks, consistent with a standard semester.

\begin{center}
\renewcommand{\arraystretch}{1.3}
\begin{longtable}{@{}c l c c@{}}
\toprule
\textbf{\#} & \textbf{Module} & \textbf{Duration} & \textbf{Syllabus \S\S} \\
\midrule
\endhead
1 & Review of Basic Properties of Functions & 2$^-$ weeks & Ch.\,1, 2.1--2.2 \\
2 & Limits and Continuity & 2 weeks & 3.1--3.5 \\
3 & Derivatives & 3$^-$ weeks & 4.1--4.11 \\
4 & Applications of Differentiation & 2$^+$ weeks & 5.1--5.5 \\
5 & Integration & 2 weeks & 5.10, 6.1--6.3 \\
6 & Applications of the Integral & 1$^+$ week & 7.1--7.3 \\
\bottomrule
\end{longtable}
\end{center}

\noindent\textbf{Overall course prerequisite (entry-level).} Solid precalculus /
college algebra \& trigonometry: function notation and graphing, exponent and
logarithm rules, trigonometric identities and the unit circle, polynomial
factoring, and algebraic manipulation of rational expressions. These are assumed
on day one and are only briefly reviewed in Module 1.

\noindent\textbf{Course-level learning arc.} Functions $\rightarrow$ limiting behavior
$\rightarrow$ instantaneous rate of change (derivatives) $\rightarrow$ using derivatives
to analyze behavior $\rightarrow$ accumulation (integrals) $\rightarrow$ techniques to
evaluate integrals. Each module is the conceptual prerequisite for the next.

% =========================================================
\section{Module 1: Review of Basic Properties of Functions}
\label{mod:1}
% =========================================================
\begin{modulebox}{Module 1 --- Functions, Exponential Growth \& Decay, Sequences \hfill (2$^-$ weeks)}

\textbf{Chapters/Sections:}
\begin{itemize}[nosep]
  \item Chapter 1 --- Preliminaries: functions, domain/range, function types, graph transformations, composition and inverses
  \item 2.1 --- Exponential Growth and Decay
  \item 2.2 --- Sequences
\end{itemize}

\textbf{Core Concepts:}
\begin{itemize}[nosep]
  \item Function definitions: polynomial, rational, power, trigonometric, exponential, logarithmic, piecewise
  \item Domain, range, and graphical behavior; even/odd symmetry; increasing/decreasing behavior
  \item Composition of functions and inverse functions
  \item Exponential models $N(t) = N_0 e^{kt}$; growth vs.\ decay; doubling time and half-life
  \item Sequences: explicit vs.\ recursive definitions; convergence, divergence, and limits of a sequence
\end{itemize}

\loheader
\begin{itemize}[nosep]
  \item Classify a function by type and determine its domain and range algebraically and graphically
  \item Compose and invert functions, and interpret inverses in applied contexts
  \item Construct exponential growth/decay models from word problems and solve for rate, time, or initial quantity
  \item Compute and interpret half-life and doubling time
  \item Determine whether a sequence converges or diverges, and find its limit when it exists
\end{itemize}
\end{modulebox}

\begin{diffbox}
\begin{itemize}[nosep]
  \item Translating a word problem into an exponential model (identifying $N_0$ and $k$ correctly, including sign of $k$)
  \item Distinguishing sequence convergence from the behavior of a related continuous function
  \item Working fluently with logarithm and exponent rules under time pressure
\end{itemize}
\end{diffbox}

\begin{prereqbox}
Algebra I/II: function notation, graphing, solving equations. Precalculus: exponent
and logarithm laws, interpreting graphs of $e^x$ and $\ln x$.
\end{prereqbox}

% =========================================================
\section{Module 2: Limits and Continuity}
\label{mod:2}
% =========================================================
\begin{modulebox}{Module 2 --- Limits, Continuity, and Limiting Behavior \hfill (2 weeks)}

\textbf{Chapters/Sections:}
\begin{itemize}[nosep]
  \item 3.1 --- Limits (intuitive approach, with the formal $\varepsilon$--$\delta$ definition introduced via 3.6)
  \item 3.2 --- Continuity
  \item 3.3 --- Limits at Infinity (with the formal definition via 3.6)
  \item 3.4 --- Trigonometric Limits and the Sandwich (Squeeze) Theorem
  \item 3.5 --- Properties of Continuous Functions
\end{itemize}

\textbf{Core Concepts:}
\begin{itemize}[nosep]
  \item One-sided and two-sided limits; limit laws
  \item Formal $\varepsilon$--$\delta$ definition of a limit
  \item Continuity at a point and on an interval; types of discontinuities (removable, jump, infinite)
  \item Horizontal/vertical asymptotes and limits at infinity
  \item Squeeze (Sandwich) Theorem, especially applied to $\lim_{\theta \to 0}\frac{\sin\theta}{\theta}$
  \item Intermediate Value Theorem (IVT) and Extreme Value Theorem groundwork
\end{itemize}

\loheader
\begin{itemize}[nosep]
  \item Evaluate limits algebraically (factoring, rationalizing), graphically, and numerically
  \item Write and interpret an $\varepsilon$--$\delta$ proof for a simple limit statement
  \item Classify discontinuities and determine continuity of piecewise functions
  \item Evaluate limits at infinity and identify horizontal/vertical asymptotes
  \item Apply the Squeeze Theorem to evaluate trigonometric limits
  \item Apply the IVT to justify the existence of roots or solutions
\end{itemize}
\end{modulebox}

\begin{diffbox}
\begin{itemize}[nosep]
  \item The formal $\varepsilon$--$\delta$ definition: translating a symbolic statement into a working proof
  \item Indeterminate forms ($0/0$, $\infty - \infty$) requiring algebraic manipulation before a limit can be evaluated
  \item Recognizing when/how to apply the Squeeze Theorem rather than direct substitution
  \item Distinguishing a removable discontinuity from a true break in continuity
\end{itemize}
\end{diffbox}

\begin{prereqbox}
Module 1 (function behavior, graphing). Algebraic fluency: factoring, rationalizing
radicals, trigonometric identities.
\end{prereqbox}

% =========================================================
\section{Module 3: Derivatives}
\label{mod:3}
% =========================================================
\begin{modulebox}{Module 3 --- Differentiation Rules and Techniques \hfill (3$^-$ weeks)}

\textbf{Chapters/Sections:}
\begin{itemize}[nosep]
  \item 4.1 --- Formal Definition of the Derivative
  \item 4.2 --- Basic Rules of Derivatives
  \item 4.3 --- Power Rule and the Rules of Differentiation
  \item 4.4 --- Product and Quotient Rules
  \item 4.5 --- The Chain Rule
  \item 4.6 --- Implicit Functions and Implicit Differentiation
  \item 4.7 --- Higher Derivatives
  \item 4.8 --- Derivatives of Trigonometric Functions
  \item 4.9 --- Derivatives of Exponential Functions
  \item 4.10 --- Derivatives of Inverse Functions
  \item 4.11 --- Linear Approximation
\end{itemize}

\textbf{Core Concepts:}
\begin{itemize}[nosep]
  \item Derivative as a limit of a difference quotient; derivative as instantaneous rate of change and slope of tangent line
  \item Differentiability vs.\ continuity
  \item Sum, constant multiple, power, product, quotient, and chain rules
  \item Implicit differentiation for curves not expressed as $y = f(x)$
  \item Second and higher-order derivatives; notation ($f''$, $\frac{d^2y}{dx^2}$)
  \item Derivatives of $\sin x$, $\cos x$, $\tan x$, $e^x$, and inverse functions (including logarithmic and inverse trig)
  \item Linear approximation / tangent line approximation, $L(x) = f(a) + f'(a)(x-a)$
\end{itemize}

\loheader
\begin{itemize}[nosep]
  \item Compute a derivative directly from the limit definition
  \item Apply the power, product, quotient, and chain rules fluently, including in combination
  \item Perform implicit differentiation and find slopes of tangent lines to implicit curves
  \item Compute higher-order derivatives and interpret their meaning (e.g., concavity, acceleration)
  \item Differentiate trigonometric, exponential, logarithmic, and inverse functions
  \item Use linear approximation to estimate function values near a point
\end{itemize}
\end{modulebox}

\begin{diffbox}
\begin{itemize}[nosep]
  \item Chain rule with nested/multiple compositions (``chain rule of chain rules'')
  \item Combining rules correctly in a single problem (e.g., quotient rule containing a chain-rule term)
  \item Implicit differentiation: correctly applying $\frac{d}{dx}$ to both sides and solving for $\frac{dy}{dx}$
  \item Derivatives of inverse functions (formula involving $1/f'(f^{-1}(x))$) and inverse trig derivatives
  \item Keeping track of higher derivatives without sign/simplification errors
\end{itemize}
\end{diffbox}

\begin{prereqbox}
Module 2 (limits, continuity). Trigonometric identities, exponential/logarithm
rules, and strong algebraic manipulation (fractions, radicals, expanding/factoring).
\end{prereqbox}

% =========================================================
\section{Module 4: Applications of Differentiation}
\label{mod:4}
% =========================================================
\begin{modulebox}{Module 4 --- Using Derivatives to Analyze Functions \hfill (2$^+$ weeks)}

\textbf{Chapters/Sections:}
\begin{itemize}[nosep]
  \item 5.1 --- Extreme and Mean Value Theorems
  \item 5.2 --- Monotonicity and Concavity
  \item 5.3 --- Extrema and Inflection Points
  \item 5.4 --- Optimization
  \item 5.5 --- L'H\^{o}pital's Rule
\end{itemize}

\textbf{Core Concepts:}
\begin{itemize}[nosep]
  \item Extreme Value Theorem (EVT) and Mean Value Theorem (MVT); Rolle's Theorem as a special case
  \item First derivative test: increasing/decreasing intervals from the sign of $f'$
  \item Second derivative test: concavity and inflection points from the sign of $f''$
  \item Local vs.\ global (absolute) extrema; critical points
  \item Constrained optimization: translating a real-world scenario into an objective function and constraint
  \item L'H\^{o}pital's Rule for indeterminate forms $\frac{0}{0}$ and $\frac{\infty}{\infty}$ (and reducible forms)
\end{itemize}

\loheader
\begin{itemize}[nosep]
  \item State and apply the EVT and MVT, and interpret their geometric meaning
  \item Find critical points and classify them using the first- and/or second-derivative tests
  \item Sketch a function's graph using information about increasing/decreasing behavior and concavity
  \item Set up and solve applied optimization problems, including biology/social-science contexts
  \item Recognize indeterminate forms and apply L'H\^{o}pital's Rule appropriately (including after algebraic rewriting)
\end{itemize}
\end{modulebox}

\begin{diffbox}
\begin{itemize}[nosep]
  \item Optimization word problems: identifying the correct objective function and constraint, and justifying a max/min with the correct test
  \item Distinguishing critical points that are extrema from those that are neither (e.g., $f(x) = x^3$ at $x=0$)
  \item Rewriting non-standard indeterminate forms ($0 \cdot \infty$, $\infty - \infty$, $0^0$, $1^\infty$) into a form usable by L'H\^{o}pital's Rule
  \item Full curve-sketching problems that require synthesizing first- and second-derivative information simultaneously
\end{itemize}
\end{diffbox}

\begin{prereqbox}
Module 3 (all differentiation rules). Comfort setting up equations from verbal
descriptions (algebra word-problem skills).
\end{prereqbox}

% =========================================================
\section{Module 5: Integration}
\label{mod:5}
% =========================================================
\begin{modulebox}{Module 5 --- Antiderivatives, the Definite Integral, and the FTC \hfill (2 weeks)}

\textbf{Chapters/Sections:}
\begin{itemize}[nosep]
  \item 5.10 --- Antiderivatives
  \item 6.1 --- The Definite Integral
  \item 6.2 --- The Fundamental Theorem of Calculus
  \item 6.3 --- Applications of Integration
\end{itemize}

\textbf{Core Concepts:}
\begin{itemize}[nosep]
  \item Antiderivatives and indefinite integrals; the constant of integration
  \item Riemann sums (left, right, midpoint) as approximations of area
  \item The definite integral as a limit of Riemann sums
  \item Fundamental Theorem of Calculus, Parts 1 and 2 --- linking differentiation and integration
  \item Applications: area between curves, average value of a function, accumulated change from a rate
\end{itemize}

\loheader
\begin{itemize}[nosep]
  \item Compute basic antiderivatives using reverse differentiation rules
  \item Set up and approximate a definite integral using Riemann sums
  \item State and apply both parts of the Fundamental Theorem of Calculus
  \item Evaluate definite integrals and interpret them as net accumulated change
  \item Apply integrals to compute area between curves and the average value of a function
\end{itemize}
\end{modulebox}

\begin{diffbox}
\begin{itemize}[nosep]
  \item Conceptually connecting Riemann sums (a discrete process) to the definite integral (a continuous limit)
  \item Correctly applying FTC Part 1 when the integrand's upper limit is itself a function (requires the chain rule)
  \item Interpreting the definite integral in applied/contextual problems (net change vs.\ total distance, signed area)
\end{itemize}
\end{diffbox}

\begin{prereqbox}
Module 3--4 (derivatives and their behavior). Summation notation and basic limit
concepts from Module 2 are used when introducing Riemann sums.
\end{prereqbox}

% =========================================================
\section{Module 6: Applications of the Integral}
\label{mod:6}
% =========================================================
\begin{modulebox}{Module 6 --- Integration Techniques \hfill (1$^+$ week)}

\textbf{Chapters/Sections:}
\begin{itemize}[nosep]
  \item 7.1 --- The Substitution Rule
  \item 7.2 --- Integration by Parts and Practicing Integration
  \item 7.3 --- Rational Functions and Partial Fractions
\end{itemize}

\textbf{Core Concepts:}
\begin{itemize}[nosep]
  \item $u$-substitution as the reverse of the chain rule; adjusting bounds for definite integrals
  \item Integration by parts, $\int u\,dv = uv - \int v\,du$, and choosing $u$ vs.\ $dv$ (e.g., via LIATE)
  \item Strategic selection between substitution, parts, and algebraic manipulation
  \item Partial fraction decomposition for rational integrands: distinct linear factors, repeated factors, irreducible quadratics
\end{itemize}

\loheader
\begin{itemize}[nosep]
  \item Apply $u$-substitution to evaluate indefinite and definite integrals, adjusting limits correctly
  \item Apply integration by parts, including cases requiring repeated application or a tabular method
  \item Decompose a rational function into partial fractions and integrate each term
  \item Select an appropriate integration technique (or combination of techniques) for a given integrand
\end{itemize}
\end{modulebox}

\begin{diffbox}
\begin{itemize}[nosep]
  \item Recognizing which technique applies (substitution vs.\ parts vs.\ partial fractions) when the problem gives no hint
  \item Integration by parts requiring two or more repeated applications (e.g., $\int x^2 e^x\,dx$) without sign errors
  \item Setting up partial fraction decomposition correctly for repeated linear factors or irreducible quadratic factors
  \item Adjusting the limits of integration correctly when using substitution on a definite integral
\end{itemize}
\end{diffbox}

\begin{prereqbox}
Module 5 (antiderivatives, definite integral). Module 3 (chain rule, product rule ---
mirrored by substitution and by-parts). Polynomial factoring and long division from
algebra.
\end{prereqbox}

% =========================================================
\section{Cross-Course Difficulty Map}
% =========================================================

The table below consolidates the topics most commonly associated with student
difficulty across the whole course, ranked roughly by where they occur.

\begin{center}
\renewcommand{\arraystretch}{1.3}
\begin{longtable}{@{}p{4.3cm}p{4.6cm}p{5.3cm}@{}}
\toprule
\textbf{Module} & \textbf{Difficult Topic} & \textbf{Why It's Difficult} \\
\midrule
\endhead
1. Functions & Building exponential models from word problems & Requires translating verbal rate information into $N_0$ and $k$ \\
2. Limits \& Continuity & $\varepsilon$--$\delta$ definition of a limit & First fully rigorous proof structure students encounter \\
2. Limits \& Continuity & Indeterminate forms before simplification & Requires algebraic insight, not just substitution \\
3. Derivatives & Chain rule with nested compositions & Multiple layers must be differentiated in the correct order \\
3. Derivatives & Implicit differentiation & Non-intuitive treatment of $y$ as a function of $x$ mid-equation \\
4. Applications & Optimization word problems & Requires modeling before any calculus is applied \\
4. Applications & Non-standard indeterminate forms ($0^0$, $1^\infty$, etc.) & Must be algebraically rewritten before L'H\^{o}pital applies \\
5. Integration & Riemann sum $\to$ definite integral connection & Abstract limiting process linking discrete sums to a continuous quantity \\
6. Integral Applications & Choosing substitution vs.\ parts vs.\ partial fractions & No single rule; requires pattern recognition built from practice \\
\bottomrule
\end{longtable}
\end{center}

% =========================================================
\section{Prerequisite Map}
% =========================================================

\begin{center}
\renewcommand{\arraystretch}{1.3}
\begin{longtable}{@{}p{4.3cm}p{9.9cm}@{}}
\toprule
\textbf{Module} & \textbf{Depends On} \\
\midrule
\endhead
1. Functions, Growth, Sequences & Algebra I/II; Precalculus (exponents, logs, graphing) \\
2. Limits \& Continuity & Module 1; algebraic factoring/rationalizing; trig identities \\
3. Derivatives & Module 2; trig identities; exponential/log rules \\
4. Applications of Differentiation & Module 3 in full \\
5. Integration & Modules 3--4; limit concepts and summation notation from Module 2 \\
6. Applications of the Integral & Module 5; chain and product rules from Module 3; polynomial factoring/division \\
\bottomrule
\end{longtable}
\end{center}

\vfill
\begin{center}
\small\textit{Compiled from the official AS.110.106 syllabus, Johns Hopkins University Department of Mathematics.}
\end{center}

\end{document}